


\documentclass[12pt]{article}
\usepackage{math_hw}
\usepackage{amssymb}


% The true underlying data generating distribution
\newcommand{\pdata}{p_{\rm{data}}}
% The empirical distribution defined by the training set
\newcommand{\ptrain}{\hat{p}_{\rm{data}}}
\newcommand{\Ptrain}{\hat{P}_{\rm{data}}}
% The model distribution
\newcommand{\pmodel}{p_{\rm{model}}}
\newcommand{\Pmodel}{P_{\rm{model}}}
\newcommand{\ptildemodel}{\tilde{p}_{\rm{model}}}
% Stochastic autoencoder distributions
\newcommand{\pencode}{p_{\rm{encoder}}}
\newcommand{\pdecode}{p_{\rm{decoder}}}
\newcommand{\precons}{p_{\rm{reconstruct}}}

\newcommand{\laplace}{\mathrm{Laplace}} % Laplace distribution

\newcommand{\E}{\mathbb{E}}
\newcommand{\Ls}{\mathcal{L}}
\newcommand{\R}{\mathbb{R}}
\newcommand{\emp}{\tilde{p}}
\newcommand{\lr}{\alpha}
\newcommand{\reg}{\lambda}
\newcommand{\rect}{\mathrm{rectifier}}
\newcommand{\softmax}{\mathrm{softmax}}
\newcommand{\sigmoid}{\sigma}
\newcommand{\softplus}{\zeta}
\newcommand{\KL}{D_{\mathrm{KL}}}
\newcommand{\Var}{\mathrm{Var}}
\newcommand{\standarderror}{\mathrm{SE}}
\newcommand{\Cov}{\mathrm{Cov}}
% Wolfram Mathworld says $L^2$ is for function spaces and $\ell^2$ is for vectors
% But then they seem to use $L^2$ for vectors throughout the site, and so does
% wikipedia.
\newcommand{\normlzero}{L^0}
\newcommand{\normlone}{L^1}
\newcommand{\normltwo}{L^2}
\newcommand{\normlp}{L^p}
\newcommand{\normmax}{L^\infty}

\newcommand{\xpos}{x_+}
\newcommand{\xneg}{x_-}

\newcommand{\parents}{Pa} % See usage in notation.tex. Chosen to match Daphne's book.

\DeclareMathOperator*{\argmax}{arg\,max}
\DeclareMathOperator*{\argmin}{arg\,min}

\DeclareMathOperator{\sign}{sign}
\let\ab\allowbreak
\newcommand{\ones}{\mathbf 1}
\newcommand{\reals}{\mathbb{R}}
\newcommand{\integers}{{\mbox{\bf Z}}}
\newcommand{\symm}{{\mbox{\bf S}}}  % symmetric matrices
\DeclareMathOperator*{\esssup}{ess\,sup}
\newcommand{\nullspace}{{\mathcal N}}
\newcommand{\range}{{\mathcal R}}
\newcommand{\diag}{\mathop{\bf diag}}
\newcommand{\card}{\mathop{\bf card}}
\newcommand{\conv}{\mathop{\bf conv}}
\newcommand{\prox}{\mathbf{Prox}}

\newcommand{\Expect}{\mathop{\bf E{}}}
\newcommand{\Prob}{\mathop{\bf Prob}}
\newcommand{\Co}{{\mathop {\bf Co}}} % convex hull
\newcommand{\dist}{\mathop{\bf dist{}}}
\newcommand{\epi}{\mathop{\bf epi}} % epigraph
\newcommand{\Vol}{\mathop{\bf vol}}
\newcommand{\dom}{\mathop{\bf dom}} % domain
\newcommand{\intr}{\mathop{\bf int}}
\newcommand{\avar}{\mathbf{AVar}}
\newcommand{\expec}{\mathbb{E}}

\newcommand{\cf}{{\it cf.}}
\newcommand{\eg}{{\it e.g.}}
\newcommand{\ie}{{\it i.e.}}
\newcommand{\etc}{{\it etc.}}


\newcommand{\fC}{\field{C}} 
\newcommand{\fE}{\field{E}} 
\newcommand{\fN}{\field{N}} % natural numbers
\newcommand{\fR}{\field{R}} % real numbers
\newcommand{\fZ}{\field{Z}} % integers

\newcommand{\bzero}[1]{0}



\newcommand{\mathify}[1]{\ifmmode{#1}\else\mbox{$#1$}\fi}
\newcommand{\Epi}[1]{\mathify{\mathbb{E}_{\pi}\left[#1\right]}}	
\newcommand{\Ep}[1]{\mathify{\mathbb{E}_{P}\left[#1\right]}}
\newcommand{\Eq}[1]{\mathify{\mathbb{E}_{Q}\left[#1\right]}}
\newcommand{\EX}[1]{\mathify{\mathbb{E}_{X}\left[#1\right]}}
\newcommand{\EZ}[1]{\mathify{\mathbb{E}_{Z}\left[#1\right]}}




\newcommand\blfootnote[1]{%
  \begingroup
  \renewcommand\thefootnote{}\footnote{#1}%
  \addtocounter{footnote}{-1}%
  \endgroup
}
%%%%% NEW MATH DEFINITIONS %%%%%

\usepackage{amsmath,amsfonts,bm}

% Mark sections of captions for referring to divisions of figures
\newcommand{\figleft}{{\em (Left)}}
\newcommand{\figcenter}{{\em (Center)}}
\newcommand{\figright}{{\em (Right)}}
\newcommand{\figtop}{{\em (Top)}}
\newcommand{\figbottom}{{\em (Bottom)}}
\newcommand{\captiona}{{\em (a)}}
\newcommand{\captionb}{{\em (b)}}
\newcommand{\captionc}{{\em (c)}}
\newcommand{\captiond}{{\em (d)}}

% Highlight a newly defined term
\newcommand{\newterm}[1]{{\bf #1}}

\newcommand{\field}[1]{\mathbb{#1}}
\newcommand{\mse}{\text{\rm mse}}
\newcommand{\Ascr}{{\cal A}}
\newcommand{\Bscr}{{\cal B}}
\newcommand{\Cscr}{{\cal C}}
\newcommand{\Dscr}{{\cal D}}
\newcommand{\Escr}{{\cal E}}
\newcommand{\Fscr}{{\cal F}}
\newcommand{\Gscr}{{\cal G}}
\newcommand{\Hscr}{{\cal H}}
\newcommand{\Iscr}{{\cal I}}
\newcommand{\Jscr}{{\cal J}}
\newcommand{\Kscr}{{\cal K}}
\newcommand{\Lscr}{{\cal L}}
\newcommand{\Mscr}{{\cal M}}
\newcommand{\Nscr}{{\cal N}}
\newcommand{\Oscr}{{\cal O}}
\newcommand{\Pscr}{{\cal P}}
\newcommand{\Qscr}{{\cal Q}}
\newcommand{\Rscr}{{\cal R}}
\newcommand{\Sscr}{{\cal S}}
\newcommand{\Tscr}{{\cal T}}
\newcommand{\Uscr}{{\cal U}}
\newcommand{\Vscr}{{\cal V}}
\newcommand{\Wscr}{{\cal W}}
\newcommand{\Xscr}{{\cal X}}
\newcommand{\Yscr}{{\cal Y}}
\newcommand{\Zscr}{{\cal Z}}

\newcommand{\Xb}[1]{{\bf X^{(#1)}}}
\newcommand{\Bb}[1]{{\bf B^{(#1)}}}
\newcommand{\Xe}[1]{ X^{(#1)}}
\newcommand{\Be}[1]{ B^{(#1)}}
\newcommand{\Zb}[1]{{\bf Z^{(#1)}}}
\newcommand{\Ze}[1]{ Z^{(#1)}}
\newcommand{\Yb}[1]{{\bf Y^{(#1)}}}
\newcommand{\Ye}[1]{ Y^{(#1)}}
\newcommand{\Ub}[1]{{\bf U^{(#1)}}}
\newcommand{\Ue}[1]{ U^{(#1)}}
\newcommand{\Vb}[1]{{\bf V^{(#1)}}}
\newcommand{\Ve}[1]{ V^{(#1)}}
\newcommand{\Wb}[1]{{\bf W^{(#1)}}}
\newcommand{\We}[1]{ W^{(#1)}}
\newcommand{\Wbtilde}[1]{{\bf \widetilde{W}^{(#1)}}}
\newcommand{\Wetilde}[1]{ \widetilde{W}^{(#1)}}
\newcommand{\Ab}{{\bf A}}
\newcommand{\Qb}{{\bf Q}}
\newcommand{\Xbhat}{{\bf \hat{X}}}
\newcommand{\Xbh}[1]{{\bf \hat{X}^{(#1)}}}
\newcommand{\mb}[1]{{\bf m^{(#1)}}}
\newcommand{\me}[1]{ m^{(#1)}}
\newcommand{\mbhat}[1]{{\bf \hat{m}^{(#1)}}}
\newcommand{\mehat}[1]{{\hat{m}^{(#1)}}}


% Figure reference, lower-case.
\def\figref#1{figure~\ref{#1}}
% Figure reference, capital. For start of sentence
\def\Figref#1{Figure~\ref{#1}}
\def\twofigref#1#2{figures \ref{#1} and \ref{#2}}
\def\quadfigref#1#2#3#4{figures \ref{#1}, \ref{#2}, \ref{#3} and \ref{#4}}
% Section reference, lower-case.
\def\secref#1{section~\ref{#1}}
% Section reference, capital.
\def\Secref#1{Section~\ref{#1}}
% Reference to two sections.
\def\twosecrefs#1#2{sections \ref{#1} and \ref{#2}}
% Reference to three sections.
\def\secrefs#1#2#3{sections \ref{#1}, \ref{#2} and \ref{#3}}
% Reference to an equation, lower-case.
\def\eqref#1{equation~\ref{#1}}
% Reference to an equation, upper case
\def\Eqref#1{Equation~\ref{#1}}
% A raw reference to an equation---avoid using if possible
\def\plaineqref#1{\ref{#1}}
% Reference to a chapter, lower-case.
\def\chapref#1{chapter~\ref{#1}}
% Reference to an equation, upper case.
\def\Chapref#1{Chapter~\ref{#1}}
% Reference to a range of chapters
\def\rangechapref#1#2{chapters\ref{#1}--\ref{#2}}
% Reference to an algorithm, lower-case.
\def\algref#1{algorithm~\ref{#1}}
% Reference to an algorithm, upper case.
\def\Algref#1{Algorithm~\ref{#1}}
\def\twoalgref#1#2{algorithms \ref{#1} and \ref{#2}}
\def\Twoalgref#1#2{Algorithms \ref{#1} and \ref{#2}}
% Reference to a part, lower case
\def\partref#1{part~\ref{#1}}
% Reference to a part, upper case
\def\Partref#1{Part~\ref{#1}}
\def\twopartref#1#2{parts \ref{#1} and \ref{#2}}

\def\ceil#1{\lceil #1 \rceil}
\def\floor#1{\lfloor #1 \rfloor}
\def\1{\bm{1}}
\newcommand{\train}{\mathcal{D}}
\newcommand{\valid}{\mathcal{D_{\mathrm{valid}}}}
\newcommand{\test}{\mathcal{D_{\mathrm{test}}}}

\def\eps{{\epsilon}}


% Random variables
\def\reta{{\textnormal{$\eta$}}}
\def\ra{{\textnormal{a}}}
\def\rb{{\textnormal{b}}}
\def\rc{{\textnormal{c}}}
\def\rd{{\textnormal{d}}}
\def\re{{\textnormal{e}}}
\def\rf{{\textnormal{f}}}
\def\rg{{\textnormal{g}}}
\def\rh{{\textnormal{h}}}
\def\ri{{\textnormal{i}}}
\def\rj{{\textnormal{j}}}
\def\rk{{\textnormal{k}}}
\def\rl{{\textnormal{l}}}
% rm is already a command, just don't name any random variables m
\def\rn{{\textnormal{n}}}
\def\ro{{\textnormal{o}}}
\def\rp{{\textnormal{p}}}
\def\rq{{\textnormal{q}}}
\def\rr{{\textnormal{r}}}
\def\rs{{\textnormal{s}}}
\def\rt{{\textnormal{t}}}
\def\ru{{\textnormal{u}}}
\def\rv{{\textnormal{v}}}
\def\rw{{\textnormal{w}}}
\def\rx{{\textnormal{x}}}
\def\ry{{\textnormal{y}}}
\def\rz{{\textnormal{z}}}

% Random vectors
\def\rvepsilon{{\mathbf{\epsilon}}}
\def\rvtheta{{\mathbf{\theta}}}
\def\rva{{\mathbf{a}}}
\def\rvb{{\mathbf{b}}}
\def\rvc{{\mathbf{c}}}
\def\rvd{{\mathbf{d}}}
\def\rve{{\mathbf{e}}}
\def\rvf{{\mathbf{f}}}
\def\rvg{{\mathbf{g}}}
\def\rvh{{\mathbf{h}}}
\def\rvu{{\mathbf{i}}}
\def\rvj{{\mathbf{j}}}
\def\rvk{{\mathbf{k}}}
\def\rvl{{\mathbf{l}}}
\def\rvm{{\mathbf{m}}}
\def\rvn{{\mathbf{n}}}
\def\rvo{{\mathbf{o}}}
\def\rvp{{\mathbf{p}}}
\def\rvq{{\mathbf{q}}}
\def\rvr{{\mathbf{r}}}
\def\rvs{{\mathbf{s}}}
\def\rvt{{\mathbf{t}}}
\def\rvu{{\mathbf{u}}}
\def\rvv{{\mathbf{v}}}
\def\rvw{{\mathbf{w}}}
\def\rvx{{\mathbf{x}}}
\def\rvy{{\mathbf{y}}}
\def\rvz{{\mathbf{z}}}

% Elements of random vectors
\def\erva{{\textnormal{a}}}
\def\ervb{{\textnormal{b}}}
\def\ervc{{\textnormal{c}}}
\def\ervd{{\textnormal{d}}}
\def\erve{{\textnormal{e}}}
\def\ervf{{\textnormal{f}}}
\def\ervg{{\textnormal{g}}}
\def\ervh{{\textnormal{h}}}
\def\ervi{{\textnormal{i}}}
\def\ervj{{\textnormal{j}}}
\def\ervk{{\textnormal{k}}}
\def\ervl{{\textnormal{l}}}
\def\ervm{{\textnormal{m}}}
\def\ervn{{\textnormal{n}}}
\def\ervo{{\textnormal{o}}}
\def\ervp{{\textnormal{p}}}
\def\ervq{{\textnormal{q}}}
\def\ervr{{\textnormal{r}}}
\def\ervs{{\textnormal{s}}}
\def\ervt{{\textnormal{t}}}
\def\ervu{{\textnormal{u}}}
\def\ervv{{\textnormal{v}}}
\def\ervw{{\textnormal{w}}}
\def\ervx{{\textnormal{x}}}
\def\ervy{{\textnormal{y}}}
\def\ervz{{\textnormal{z}}}

% Random matrices
\def\rmA{{\mathbf{A}}}
\def\rmB{{\mathbf{B}}}
\def\rmC{{\mathbf{C}}}
\def\rmD{{\mathbf{D}}}
\def\rmE{{\mathbf{E}}}
\def\rmF{{\mathbf{F}}}
\def\rmG{{\mathbf{G}}}
\def\rmH{{\mathbf{H}}}
\def\rmI{{\mathbf{I}}}
\def\rmJ{{\mathbf{J}}}
\def\rmK{{\mathbf{K}}}
\def\rmL{{\mathbf{L}}}
\def\rmM{{\mathbf{M}}}
\def\rmN{{\mathbf{N}}}
\def\rmO{{\mathbf{O}}}
\def\rmP{{\mathbf{P}}}
\def\rmQ{{\mathbf{Q}}}
\def\rmR{{\mathbf{R}}}
\def\rmS{{\mathbf{S}}}
\def\rmT{{\mathbf{T}}}
\def\rmU{{\mathbf{U}}}
\def\rmV{{\mathbf{V}}}
\def\rmW{{\mathbf{W}}}
\def\rmX{{\mathbf{X}}}
\def\rmY{{\mathbf{Y}}}
\def\rmZ{{\mathbf{Z}}}

% Elements of random matrices
\def\ermA{{\textnormal{A}}}
\def\ermB{{\textnormal{B}}}
\def\ermC{{\textnormal{C}}}
\def\ermD{{\textnormal{D}}}
\def\ermE{{\textnormal{E}}}
\def\ermF{{\textnormal{F}}}
\def\ermG{{\textnormal{G}}}
\def\ermH{{\textnormal{H}}}
\def\ermI{{\textnormal{I}}}
\def\ermJ{{\textnormal{J}}}
\def\ermK{{\textnormal{K}}}
\def\ermL{{\textnormal{L}}}
\def\ermM{{\textnormal{M}}}
\def\ermN{{\textnormal{N}}}
\def\ermO{{\textnormal{O}}}
\def\ermP{{\textnormal{P}}}
\def\ermQ{{\textnormal{Q}}}
\def\ermR{{\textnormal{R}}}
\def\ermS{{\textnormal{S}}}
\def\ermT{{\textnormal{T}}}
\def\ermU{{\textnormal{U}}}
\def\ermV{{\textnormal{V}}}
\def\ermW{{\textnormal{W}}}
\def\ermX{{\textnormal{X}}}
\def\ermY{{\textnormal{Y}}}
\def\ermZ{{\textnormal{Z}}}

% Vectors
\def\vzero{{\bm{0}}}
\def\vone{{\bm{1}}}
\def\vmu{{\bm{\mu}}}
\def\vtheta{{\bm{\theta}}}
\def\va{{\bm{a}}}
\def\vb{{\bm{b}}}
\def\vc{{\bm{c}}}
\def\vd{{\bm{d}}}
\def\ve{{\bm{e}}}
\def\vf{{\bm{f}}}
\def\vg{{\bm{g}}}
\def\vh{{\bm{h}}}
\def\vi{{\bm{i}}}
\def\vj{{\bm{j}}}
\def\vk{{\bm{k}}}
\def\vl{{\bm{l}}}
\def\vm{{\bm{m}}}
\def\vn{{\bm{n}}}
\def\vo{{\bm{o}}}
\def\vp{{\bm{p}}}
\def\vq{{\bm{q}}}
\def\vr{{\bm{r}}}
\def\vs{{\bm{s}}}
\def\vt{{\bm{t}}}
\def\vu{{\bm{u}}}
\def\vv{{\bm{v}}}
\def\vw{{\bm{w}}}
\def\vx{{\bm{x}}}
\def\vy{{\bm{y}}}
\def\vz{{\bm{z}}}

% Elements of vectors
\def\evalpha{{\alpha}}
\def\evbeta{{\beta}}
\def\evepsilon{{\epsilon}}
\def\evlambda{{\lambda}}
\def\evomega{{\omega}}
\def\evmu{{\mu}}
\def\evpsi{{\psi}}
\def\evsigma{{\sigma}}
\def\evtheta{{\theta}}
\def\eva{{a}}
\def\evb{{b}}
\def\evc{{c}}
\def\evd{{d}}
\def\eve{{e}}
\def\evf{{f}}
\def\evg{{g}}
\def\evh{{h}}
\def\evi{{i}}
\def\evj{{j}}
\def\evk{{k}}
\def\evl{{l}}
\def\evm{{m}}
\def\evn{{n}}
\def\evo{{o}}
\def\evp{{p}}
\def\evq{{q}}
\def\evr{{r}}
\def\evs{{s}}
\def\evt{{t}}
\def\evu{{u}}
\def\evv{{v}}
\def\evw{{w}}
\def\evx{{x}}
\def\evy{{y}}
\def\evz{{z}}

% Matrix
\def\mA{{\bm{A}}}
\def\mB{{\bm{B}}}
\def\mC{{\bm{C}}}
\def\mD{{\bm{D}}}
\def\mE{{\bm{E}}}
\def\mF{{\bm{F}}}
\def\mG{{\bm{G}}}
\def\mH{{\bm{H}}}
\def\mI{{\bm{I}}}
\def\mJ{{\bm{J}}}
\def\mK{{\bm{K}}}
\def\mL{{\bm{L}}}
\def\mM{{\bm{M}}}
\def\mN{{\bm{N}}}
\def\mO{{\bm{O}}}
\def\mP{{\bm{P}}}
\def\mQ{{\bm{Q}}}
\def\mR{{\bm{R}}}
\def\mS{{\bm{S}}}
\def\mT{{\bm{T}}}
\def\mU{{\bm{U}}}
\def\mV{{\bm{V}}}
\def\mW{{\bm{W}}}
\def\mX{{\bm{X}}}
\def\mY{{\bm{Y}}}
\def\mZ{{\bm{Z}}}
\def\mBeta{{\bm{\beta}}}
\def\mPhi{{\bm{\Phi}}}
\def\mLambda{{\bm{\Lambda}}}
\def\mSigma{{\bm{\Sigma}}}

% Tensor
\DeclareMathAlphabet{\mathsfit}{\encodingdefault}{\sfdefault}{m}{sl}
\SetMathAlphabet{\mathsfit}{bold}{\encodingdefault}{\sfdefault}{bx}{n}
\newcommand{\tens}[1]{\bm{\mathsfit{#1}}}
\def\tA{{\tens{A}}}
\def\tB{{\tens{B}}}
\def\tC{{\tens{C}}}
\def\tD{{\tens{D}}}
\def\tE{{\tens{E}}}
\def\tF{{\tens{F}}}
\def\tG{{\tens{G}}}
\def\tH{{\tens{H}}}
\def\tI{{\tens{I}}}
\def\tJ{{\tens{J}}}
\def\tK{{\tens{K}}}
\def\tL{{\tens{L}}}
\def\tM{{\tens{M}}}
\def\tN{{\tens{N}}}
\def\tO{{\tens{O}}}
\def\tP{{\tens{P}}}
\def\tQ{{\tens{Q}}}
\def\tR{{\tens{R}}}
\def\tS{{\tens{S}}}
\def\tT{{\tens{T}}}
\def\tU{{\tens{U}}}
\def\tV{{\tens{V}}}
\def\tW{{\tens{W}}}
\def\tX{{\tens{X}}}
\def\tY{{\tens{Y}}}
\def\tZ{{\tens{Z}}}


\newcommand{\cA}{{\cal A}}
\newcommand{\cB}{{\cal B}}
\newcommand{\cC}{{\cal C}}
\newcommand{\cD}{{\cal D}}
\newcommand{\cE}{{\cal E}}
\newcommand{\cF}{{\cal F}}
\newcommand{\cG}{{\cal G}}
\newcommand{\cH}{{\cal H}}
\newcommand{\cI}{{\cal I}}
\newcommand{\cJ}{{\cal J}}
\newcommand{\cK}{{\cal K}}
\newcommand{\cL}{{\cal L}}
\newcommand{\cM}{{\cal M}}
\newcommand{\cN}{{\cal N}}
\newcommand{\cO}{{\cal O}}
\newcommand{\cP}{{\cal P}}
\newcommand{\cQ}{{\cal Q}}
\newcommand{\cR}{{\cal R}}
\newcommand{\cS}{{\cal S}}
\newcommand{\cT}{{\cal T}}
\newcommand{\cU}{{\cal U}}
\newcommand{\cV}{{\cal V}}
\newcommand{\cW}{{\cal W}}
\newcommand{\cX}{{\cal X}}
\newcommand{\cY}{{\cal Y}}
\newcommand{\cZ}{{\cal Z}}

% Graph
\def\gA{{\mathcal{A}}}
\def\gB{{\mathcal{B}}}
\def\gC{{\mathcal{C}}}
\def\gD{{\mathcal{D}}}
\def\gE{{\mathcal{E}}}
\def\gF{{\mathcal{F}}}
\def\gG{{\mathcal{G}}}
\def\gH{{\mathcal{H}}}
\def\gI{{\mathcal{I}}}
\def\gJ{{\mathcal{J}}}
\def\gK{{\mathcal{K}}}
\def\gL{{\mathcal{L}}}
\def\gM{{\mathcal{M}}}
\def\gN{{\mathcal{N}}}
\def\gO{{\mathcal{O}}}
\def\gP{{\mathcal{P}}}
\def\gQ{{\mathcal{Q}}}
\def\gR{{\mathcal{R}}}
\def\gS{{\mathcal{S}}}
\def\gT{{\mathcal{T}}}
\def\gU{{\mathcal{U}}}
\def\gV{{\mathcal{V}}}
\def\gW{{\mathcal{W}}}
\def\gX{{\mathcal{X}}}
\def\gY{{\mathcal{Y}}}
\def\gZ{{\mathcal{Z}}}

% Sets
\def\sA{{\mathbb{A}}}
\def\sB{{\mathbb{B}}}
\def\sC{{\mathbb{C}}}
\def\sD{{\mathbb{D}}}
% Don't use a set called E, because this would be the same as our symbol
% for expectation.
\def\sF{{\mathbb{F}}}
\def\sG{{\mathbb{G}}}
\def\sH{{\mathbb{H}}}
\def\sI{{\mathbb{I}}}
\def\sJ{{\mathbb{J}}}
\def\sK{{\mathbb{K}}}
\def\sL{{\mathbb{L}}}
\def\sM{{\mathbb{M}}}
\def\sN{{\mathbb{N}}}
\def\sO{{\mathbb{O}}}
\def\sP{{\mathbb{P}}}
\def\sQ{{\mathbb{Q}}}
\def\sR{{\mathbb{R}}}
\def\sS{{\mathbb{S}}}
\def\sT{{\mathbb{T}}}
\def\sU{{\mathbb{U}}}
\def\sV{{\mathbb{V}}}
\def\sW{{\mathbb{W}}}
\def\sX{{\mathbb{X}}}
\def\sY{{\mathbb{Y}}}
\def\sZ{{\mathbb{Z}}}


% Entries of a matrix
\def\emLambda{{\Lambda}}
\def\emA{{A}}
\def\emB{{B}}
\def\emC{{C}}
\def\emD{{D}}
\def\emE{{E}}
\def\emF{{F}}
\def\emG{{G}}
\def\emH{{H}}
\def\emI{{I}}
\def\emJ{{J}}
\def\emK{{K}}
\def\emL{{L}}
\def\emM{{M}}
\def\emN{{N}}
\def\emO{{O}}
\def\emP{{P}}
\def\emQ{{Q}}
\def\emR{{R}}
\def\emS{{S}}
\def\emT{{T}}
\def\emU{{U}}
\def\emV{{V}}
\def\emW{{W}}
\def\emX{{X}}
\def\emY{{Y}}
\def\emZ{{Z}}
\def\emSigma{{\Sigma}}

% entries of a tensor
% Same font as tensor, without \bm wrapper
\newcommand{\etens}[1]{\mathsfit{#1}}
\def\etLambda{{\etens{\Lambda}}}
\def\etA{{\etens{A}}}
\def\etB{{\etens{B}}}
\def\etC{{\etens{C}}}
\def\etD{{\etens{D}}}
\def\etE{{\etens{E}}}
\def\etF{{\etens{F}}}
\def\etG{{\etens{G}}}
\def\etH{{\etens{H}}}
\def\etI{{\etens{I}}}
\def\etJ{{\etens{J}}}
\def\etK{{\etens{K}}}
\def\etL{{\etens{L}}}
\def\etM{{\etens{M}}}
\def\etN{{\etens{N}}}
\def\etO{{\etens{O}}}
\def\etP{{\etens{P}}}
\def\etQ{{\etens{Q}}}
\def\etR{{\etens{R}}}
\def\etS{{\etens{S}}}
\def\etT{{\etens{T}}}
\def\etU{{\etens{U}}}
\def\etV{{\etens{V}}}
\def\etW{{\etens{W}}}
\def\etX{{\etens{X}}}
\def\etY{{\etens{Y}}}
\def\etZ{{\etens{Z}}}


\newcommand{\Exp}{\mathrm{E}}
\newcommand{\reals}{\fR}

\newcommand{\field}[1]{\mathbb{#1}}

\excludeversion{problem}
\includeversion{solution}

\begin{document}

\begin{center}
{\bf Data Privacy: Final\begin{solution}Solutions\end{solution}} \\ \medskip
\begin{problem}June 11, 2026 \end{problem}
\end{center}

\begin{problem}
\noindent 
% Note: Submit homework via LearnUS.
% Also, we have some extra problems in Problem 7 (3-5).
\LeftFoot{Final}
\end{problem}

\begin{solution}
\LeftFoot{Final Solution}
\end{solution}

\begin{problem}

    \vskip .3in
    {\bf ID : \underline{\hspace{6cm}}}\\
    \vskip .3in
    {\bf Name : \underline{\hspace{5.5cm}}}\\
    ~
    \vskip 2in
    \begin{table}[h]
        \Large
        \begin{tabular}{c|c}
           1 &\textcolor{white}{aaaaaaaa}/~30\\
           2 &\textcolor{white}{aaaaaaaa}/~30 \\
           3 &\textcolor{white}{aaaaaaaa}/~40 \\
            \hline
           Total &  \textcolor{white}{aaaaaaaa}/~100 \\
        \end{tabular}
    \end{table}
    \newpage
    {\bf extra sheet}
    \newpage
\end{problem}



\begin{enumerate}

%%%%%%%%%%%%%%%%%%%%%%%%%%%%% Problem 5 %%%%%%%%%%%%%%%%%%%%%%%%%%%%%
\item \textbf{LiRA Closed Form and Decision Boundary.}
Assume that for a target record $z$, the per-shadow-model loss has Gaussian distributions under the two membership hypotheses:
\begin{align*}
    \ell \mid (z \in D) \sim \cN(\mu_{\mathrm{in}}, \sigma_{\mathrm{in}}^2), \qquad \ell \mid (z \notin D) \sim \cN(\mu_{\mathrm{out}}, \sigma_{\mathrm{out}}^2).
\end{align*}
\begin{enumerate}
    \item Derive the log-likelihood ratio
    \begin{align*}
        \log \Lambda(\ell) = \log \frac{\cN(\ell ; \mu_{\mathrm{in}}, \sigma_{\mathrm{in}}^2)}{\cN(\ell ; \mu_{\mathrm{out}}, \sigma_{\mathrm{out}}^2)} = \frac{(\ell - \mu_{\mathrm{out}})^2}{2\sigma_{\mathrm{out}}^2} - \frac{(\ell - \mu_{\mathrm{in}})^2}{2\sigma_{\mathrm{in}}^2} + \log\frac{\sigma_{\mathrm{out}}}{\sigma_{\mathrm{in}}}.
    \end{align*}

    \item Show that for general $\sigma_{\mathrm{in}} \ne \sigma_{\mathrm{out}}$, the optimal decision region $\{\log \Lambda > \log \tau\}$ is one of: a bounded interval, the complement of a bounded interval, a half-line, all of $\fR$, or empty. Identify which case occurs based on the sign of $\sigma_{\mathrm{in}} - \sigma_{\mathrm{out}}$.

    \item Show that in the homoscedastic case $\sigma_{\mathrm{in}} = \sigma_{\mathrm{out}} = \sigma$, the optimal MIA reduces to the single threshold $\cA(z) = \mathbf{1}[\ell \le c]$, and express $c$ in terms of $\mu_{\mathrm{in}}, \mu_{\mathrm{out}}, \sigma, \tau$.

    \item In the homoscedastic case, derive the closed-form ROC: at false-positive rate $\alpha$, the optimal true-positive rate is
    \begin{align*}
        \mathrm{TPR}(\alpha) = \Phi\!\left(\Phi^{-1}(\alpha) + \frac{\mu_{\mathrm{out}} - \mu_{\mathrm{in}}}{\sigma}\right),
    \end{align*}
    where $\Phi$ is the standard normal CDF.
\end{enumerate}

%%%%%%%%%%%%%%%%%%%%%%%% Problem 5 Solution %%%%%%%%%%%%%%%%%%%%%%%%%%
\begin{solution}
{\bf Solution:}\\
\begin{enumerate}
    \item The Gaussian log-density is $\log\cN(\ell; \mu, \sigma^2) = -\frac{(\ell - \mu)^2}{2\sigma^2} - \frac{1}{2}\log(2\pi\sigma^2)$. Subtracting,
    \begin{align*}
        \log\Lambda(\ell) = \frac{(\ell - \mu_{\mathrm{out}})^2}{2\sigma_{\mathrm{out}}^2} - \frac{(\ell - \mu_{\mathrm{in}})^2}{2\sigma_{\mathrm{in}}^2} + \log\frac{\sigma_{\mathrm{out}}}{\sigma_{\mathrm{in}}}.
    \end{align*}

    \item Expanding, $\log\Lambda(\ell)$ is a quadratic in $\ell$ with leading coefficient
    \begin{align*}
        A = \frac{1}{2}\left(\frac{1}{\sigma_{\mathrm{out}}^2} - \frac{1}{\sigma_{\mathrm{in}}^2}\right).
    \end{align*}
    \begin{itemize}
        \item $\sigma_{\mathrm{in}} < \sigma_{\mathrm{out}}$ (members are more concentrated): $A < 0$, so $\log\Lambda$ is concave in $\ell$. Then $\{\log\Lambda > \log\tau\}$ is a (possibly empty, possibly unbounded) \emph{interval} $(c_-, c_+)$ on the loss axis.
        \item $\sigma_{\mathrm{in}} > \sigma_{\mathrm{out}}$: $A > 0$, so $\log\Lambda$ is convex. Then $\{\log\Lambda > \log\tau\}$ is the \emph{complement of a bounded interval}: $\{\ell < c_-\} \cup \{\ell > c_+\}$.
        \item $\sigma_{\mathrm{in}} = \sigma_{\mathrm{out}}$: $A = 0$, $\log\Lambda$ is linear, and the region is a half-line.
    \end{itemize}

    \item \textbf{Homoscedastic.} With $\sigma_{\mathrm{in}} = \sigma_{\mathrm{out}} = \sigma$ the quadratic terms cancel:
    \begin{align*}
        \log\Lambda(\ell) = \frac{(\ell - \mu_{\mathrm{out}})^2 - (\ell - \mu_{\mathrm{in}})^2}{2\sigma^2} = \frac{(\mu_{\mathrm{out}} - \mu_{\mathrm{in}})(\mu_{\mathrm{in}} + \mu_{\mathrm{out}} - 2\ell)}{2\sigma^2},
    \end{align*}
    a strictly decreasing affine function of $\ell$ (since $\mu_{\mathrm{out}} > \mu_{\mathrm{in}}$). Hence $\log\Lambda \ge \log\tau$ iff $\ell \le c$, where
    \begin{align*}
        c = \frac{\mu_{\mathrm{in}} + \mu_{\mathrm{out}}}{2} - \frac{\sigma^2 \log\tau}{\mu_{\mathrm{out}} - \mu_{\mathrm{in}}}.
    \end{align*}

    \item $\mathrm{FPR}(c) = \Pr_{\mathrm{out}}[\ell \le c] = \Phi((c - \mu_{\mathrm{out}})/\sigma)$. Solving for $c$ at $\mathrm{FPR} = \alpha$: $c = \mu_{\mathrm{out}} + \sigma\Phi^{-1}(\alpha)$. Substituting,
    \begin{align*}
        \mathrm{TPR}(\alpha) = \Pr_{\mathrm{in}}[\ell \le c] = \Phi\!\left(\frac{c - \mu_{\mathrm{in}}}{\sigma}\right) = \Phi\!\left(\Phi^{-1}(\alpha) + \frac{\mu_{\mathrm{out}} - \mu_{\mathrm{in}}}{\sigma}\right).
    \end{align*}
    The discrimination is controlled entirely by the standardized gap $(\mu_{\mathrm{out}} - \mu_{\mathrm{in}})/\sigma$.
\end{enumerate}
\end{solution}

\vskip .2in
\vskip .2in

%%%%%%%%%%%%%%%%%%%%%%%%%%%%% Problem 10 %%%%%%%%%%%%%%%%%%%%%%%%%%%%%
\item \textbf{Black-Box Unlearning Audits Require Many Queries.}
Let $\theta_o$ and $\theta_u$ be two LLMs with output distributions $P_0$ and $P_1$ on a fixed prompt distribution. For two distributions $P, Q$ on the same sample space, define the \emph{total variation distance}
\begin{align*}
    \mathrm{TV}(P, Q) = \sup_A |P(A) - Q(A)|,
\end{align*}
where the supremum is over events $A$. Assume $\mathrm{TV}(P_0, P_1) \le \eta$. An auditor draws $q$ i.i.d.\ prompts $x_1, \ldots, x_q$, observes one sample $y_i \sim P_b(\cdot \mid x_i)$ per prompt, and outputs $\hat b \in \{0, 1\}$. Write $P_b^{\otimes q}$ for the product distribution of the $q$ observations under $b$.
\begin{enumerate}
    \item \textbf{(Le Cam two-point lower bound.)} Show that for any test based on $q$ samples,
    \begin{align*}
        \Pr_{b=0}[\hat b \ne 0] + \Pr_{b=1}[\hat b \ne 1] \ge 1 - \mathrm{TV}(P_0^{\otimes q}, P_1^{\otimes q}).
    \end{align*}
    Hint: Identify $\hat b$ with a function $\psi : \cX^q \to [0, 1]$ (the probability of declaring $b = 1$). Then prove the auxiliary fact: for any distributions $P, Q$ and any $\psi: \cX^q \to [0, 1]$, $\fE_P[\psi] - \fE_Q[\psi] \le \mathrm{TV}(P, Q)$.\\
    (For the auxiliary fact, write $\psi(x) = \int_0^1 \mathbf{1}[\psi(x) > s]\, ds$ and integrate.)

    \item Show that $\mathrm{TV}(P_0^{\otimes q}, P_1^{\otimes q}) \le q \cdot \mathrm{TV}(P_0, P_1) \le q\eta$, by a hybrid (telescoping) argument: define $H_k = P_0^{\otimes k} \otimes P_1^{\otimes (q-k)}$ for $k = 0, \ldots, q$, and bound $\mathrm{TV}(H_k, H_{k-1})$.

    \item Conclude that to guarantee total error $\le 2\alpha$ (i.e., $\Pr_0[\hat b \ne 0] + \Pr_1[\hat b \ne 1] \le 2\alpha$), the auditor needs
    \begin{align*}
        q \ge \frac{1 - 2\alpha}{\eta}.
    \end{align*}
    Tighter $\chi^2$-style bounds give $q = \Omega(1/\eta^2)$. Why does this make single-prompt audits of LLM unlearning fundamentally insufficient when $\theta_u, \theta_o$ differ only by a small distributional perturbation?
\end{enumerate}

%%%%%%%%%%%%%%%%%%%%%%%% Problem 10 Solution %%%%%%%%%%%%%%%%%%%%%%%%%%
\begin{solution}
{\bf Solution:}\\
\begin{enumerate}
    \item \textbf{Auxiliary fact.} For any $\psi : \cX^q \to [0, 1]$, write $\psi(x) = \int_0^1 \mathbf{1}[\psi(x) > s]\, ds$ (since $\int_0^1 \mathbf{1}[\psi(x) > s] ds$ equals the length of $\{s \in [0, 1] : s < \psi(x)\} = [0, \psi(x))$, which is $\psi(x)$). Integrating against any distribution $R$ and swapping the order,
    \begin{align*}
        \fE_R[\psi] = \int_0^1 R(\{\psi > s\})\, ds.
    \end{align*}
    Therefore
    \begin{align*}
        \fE_P[\psi] - \fE_Q[\psi] = \int_0^1 [P(\{\psi > s\}) - Q(\{\psi > s\})]\, ds \le \int_0^1 \mathrm{TV}(P, Q)\, ds = \mathrm{TV}(P, Q),
    \end{align*}
    where the inequality uses $P(A) - Q(A) \le \mathrm{TV}(P, Q)$ for every event $A$.

    \textbf{Le Cam.} Identify $\hat b$ with the test $\psi(\cdot) = \Pr[\hat b = 1 \mid \cdot] \in [0, 1]$. With $P = P_1^{\otimes q}$ and $Q = P_0^{\otimes q}$,
    \begin{align*}
        \Pr_0[\hat b \ne 0] + \Pr_1[\hat b \ne 1] = \fE_Q[\psi] + \fE_P[1 - \psi] = 1 - (\fE_P[\psi] - \fE_Q[\psi]) \ge 1 - \mathrm{TV}(P, Q).
    \end{align*}

    \item \textbf{Hybrid argument.} Set $H_k = P_0^{\otimes k} \otimes P_1^{\otimes (q-k)}$, so $H_q = P_0^{\otimes q}$ and $H_0 = P_1^{\otimes q}$. The TV distance satisfies the triangle inequality (apply $|P(A) - Q(A)| \le |P(A) - R(A)| + |R(A) - Q(A)|$ and take supremum over $A$), so
    \begin{align*}
        \mathrm{TV}(P_0^{\otimes q}, P_1^{\otimes q}) = \mathrm{TV}(H_q, H_0) \le \sum_{k=1}^q \mathrm{TV}(H_k, H_{k-1}).
    \end{align*}
    Each pair $H_k, H_{k-1}$ differs only in coordinate $k$: in $H_k$ the $k$-th coordinate is drawn from $P_0$, in $H_{k-1}$ from $P_1$. For any event $A \subseteq \cX^q$, fix the other coordinates $x_{-k}$ and let $A_{x_{-k}} = \{y : (x_1, \ldots, x_{k-1}, y, x_{k+1}, \ldots, x_q) \in A\}$ be the section. Then
    \begin{align*}
        H_k(A) - H_{k-1}(A) = \fE_{x_{-k}}\!\left[P_0(A_{x_{-k}}) - P_1(A_{x_{-k}})\right],
    \end{align*}
    and the inner difference is $\le \mathrm{TV}(P_0, P_1)$ for every $x_{-k}$. Hence $H_k(A) - H_{k-1}(A) \le \mathrm{TV}(P_0, P_1)$; taking the supremum over $A$ gives $\mathrm{TV}(H_k, H_{k-1}) \le \mathrm{TV}(P_0, P_1) \le \eta$. Summing the telescope, $\mathrm{TV}(P_0^{\otimes q}, P_1^{\otimes q}) \le q\eta$.

    \item Combining (a) and (b),
    \begin{align*}
        2\alpha \ge \Pr_0[\hat b \ne 0] + \Pr_1[\hat b \ne 1] \ge 1 - q\eta \implies q \ge \frac{1 - 2\alpha}{\eta}.
    \end{align*}
    For $\alpha = 0.1$ this needs $q \ge 0.8/\eta$. Tighter Pinsker/Hellinger bounds with $\mathrm{TV}(P^{\otimes q}, Q^{\otimes q})^2 \le 2q\,\mathrm{KL}(P\|Q)$ give the $q = \Omega(1/\eta^2)$ scaling (matching the $\chi^2$ form behind Pawelczyk et al.\ 2024).

    \textbf{Implication for unlearning.} If $\theta_u$ and $\theta_o$ differ in output distribution by a small TV $\eta$, then a black-box auditor with fewer than $\approx 1/\eta$ queries cannot reliably distinguish them. Single-prompt audits --- the common regulatory-checkbox approach --- fundamentally cannot certify that unlearning has happened; the auditor must use distributional probes (many queries) to detect the perturbation.
\end{enumerate}
\end{solution}
%%%%%%%%%%%%%%%%%%%%%%%%%%%%%%%%%%%%%%%%%%%%%%%%%%%%%%%%%%%
%%% Problem
%%%%%%%%%%%%%%%%%%%%%%%%%%%%%%%%%%%%%%%%%%%%%%%%%%%%%%%%%%%
\item {\bf Nonuniform Time Scale Diffusion} {(\it 30 points)}\\

\vskip .2in
\begin{solution}{\bf Solution: Nonuniform Time Scale Diffusion.}
\begin{enumerate}
    \item Since
\end{enumerate}
\end{solution}



\begin{problem}
    \newpage~\newpage
\end{problem}


%%%%%%%%%%%%%%%%%%%%%%%%%%%%%%%%%%%%%%%%%%%%%%%%%%%%%%%%%%%
%%% Problem
%%%%%%%%%%%%%%%%%%%%%%%%%%%%%%%%%%%%%%%%%%%%%%%%%%%%%%%%%%%
\vskip .3in
\item {\bf Masked Diffusion} {\it (30 points)}\\
N ote that

\vskip .2in
\begin{solution}{\bf Solution: Masked Diffusion.}
\end{solution}




\begin{problem}
    \newpage~\newpage
\end{problem}



%%%%%%%%%%%%%%%%%%%%%%%%%%%%%%%%%%%%%%%%%%%%%%%%%%%%%%%%%%%
%%% Problem
%%%%%%%%%%%%%%%%%%%%%%%%%%%%%%%%%%%%%%%%%%%%%%%%%%%%%%%%%%%
\vskip .3in
\item {\bf Discrete Langevin} {\it (40 points)}\\

\vskip .2in
\begin{solution}{\bf Solution: Discrete Langevin.}
\end{solution}


\begin{problem}
    \newpage~\newpage~\newpage~\newpage
\end{problem}

\end{enumerate}



\end{document}
%%% Problem for 2026
Do the similar for $X_0\sim \cN(\mu, \sigma^2)$
Compute the exact score function of DDPM for all $X_t$.
Compute the distribution of $\tilde{X}_0$ when we sample from $X_n\sim\mathcal{N}(0, 1)$.
See the difference between $p_{X_0}$ and $p_{\tilde{X}_0}$.

